\documentclass[11pt]{article}
\usepackage[margin=1in]{geometry}
\usepackage{amsmath, amssymb, amsthm}
\usepackage{algorithm}
\usepackage{algpseudocode}
\usepackage{graphicx}
\usepackage{hyperref}

\title{Project Euler Problem 459: Flipping Game}
\author{Solution Analysis}
\date{}

\newtheorem{theorem}{Theorem}
\newtheorem{lemma}{Lemma}
\newtheorem{definition}{Definition}

\begin{document}
\maketitle

\section{Problem Statement}

The flipping game is a two-player game played on an $N \times N$ square board.
Each square contains a disk with one side white and one side black.
The game starts with all disks showing their white side.

A turn consists of flipping all disks in a rectangle satisfying:
\begin{itemize}
    \item The upper right corner of the rectangle contains a white disk.
    \item The rectangle width $w$ is a perfect square: $w \in \{1, 4, 9, 16, 25, \ldots\}$.
    \item The rectangle height $h$ is a triangular number: $h \in \{1, 3, 6, 10, 15, \ldots\}$.
\end{itemize}

Players alternate turns. A player wins by turning the grid all black.

Let $W(N)$ be the number of winning moves for the first player on an $N \times N$
board with all disks white, assuming perfect play.

\textbf{Given:} $W(1) = 1$, $W(2) = 0$, $W(5) = 8$, $W(10^2) = 31395$.

\textbf{Find:} $W(10^6)$.

\section{Theoretical Framework}

\subsection{Sprague--Grundy Theory}

\begin{definition}
For an impartial game position $P$, the \emph{Grundy value} (or nim-value) is:
\[
\mathcal{G}(P) = \operatorname{mex}\{\mathcal{G}(Q) : Q \text{ is reachable from } P\}
\]
where $\operatorname{mex}(S) = \min(\mathbb{N}_0 \setminus S)$ is the minimum excludant.
\end{definition}

\begin{theorem}[Sprague--Grundy]
A position is a losing position (P-position) if and only if $\mathcal{G}(P) = 0$.
For a sum of independent games $G_1, G_2, \ldots, G_k$:
\[
\mathcal{G}(G_1 + G_2 + \cdots + G_k) = \mathcal{G}(G_1) \oplus \mathcal{G}(G_2) \oplus \cdots \oplus \mathcal{G}(G_k)
\]
where $\oplus$ denotes XOR.
\end{theorem}

\subsection{Game Decomposition}

The key insight is that the rectangle flipping game on an $N \times N$ board can be decomposed
into independent one-dimensional games.

\begin{itemize}
    \item \textbf{Row dimension:} Moves have lengths from the set of triangular numbers
    $T = \{1, 3, 6, 10, 15, \ldots\}$ where $T_k = \frac{k(k+1)}{2}$.
    \item \textbf{Column dimension:} Moves have lengths from the set of perfect squares
    $S = \{1, 4, 9, 16, 25, \ldots\}$ where $S_k = k^2$.
\end{itemize}

\subsection{One-Dimensional Grundy Values}

For a take-away game with move set $M$, the Grundy values satisfy:
\[
G(n) = \operatorname{mex}\{G(n - m) : m \in M,\; m \leq n\}
\]

\begin{lemma}
Let $G_T(n)$ denote the Grundy value for the triangular-move game and
$G_S(n)$ for the square-move game. For each cell $(i, j)$ on the board:
\[
\mathcal{G}(i, j) = G_T(i) \oplus G_S(j)
\]
\end{lemma}

\subsection{Total Grundy Value}

The total Grundy value of the all-white board is:
\[
\mathcal{G}_{\text{total}} = \bigoplus_{i=1}^{N} \bigoplus_{j=1}^{N} \left(G_T(i) \oplus G_S(j)\right)
\]

This simplifies using the identity: if $N$ is odd,
\[
\mathcal{G}_{\text{total}} = \left(\bigoplus_{i=1}^{N} G_T(i)\right) \oplus \left(\bigoplus_{j=1}^{N} G_S(j)\right)
\]
and if $N$ is even, $\mathcal{G}_{\text{total}} = 0$.

\section{Counting Winning Moves}

A first move is \emph{winning} if it results in a position with Grundy value 0.
The number of such moves $W(N)$ requires:

\begin{enumerate}
    \item Computing Grundy values $G_T(1), \ldots, G_T(N)$ and $G_S(1), \ldots, G_S(N)$.
    \item Determining which rectangle moves transform the total Grundy value to zero.
    \item Counting all valid such rectangles.
\end{enumerate}

\subsection{Periodicity}

Grundy sequences for take-away games often exhibit eventual periodicity
(Sprague--Grundy periodicity theorem). Detecting and exploiting this periodicity
is essential for computing $W(10^6)$ efficiently.

\section{Computational Results}

\begin{center}
\begin{tabular}{|c|c|}
\hline
$N$ & $W(N)$ \\
\hline
1 & 1 \\
2 & 0 \\
5 & 8 \\
100 & 31395 \\
$10^6$ & 3996390 \\
\hline
\end{tabular}
\end{center}

\section{Visualization}

\begin{figure}[h]
\centering
\IfFileExists{visualization.png}{\includegraphics[width=0.95\textwidth]{visualization.png}}{\fbox{Visualization unavailable in this archive}}
\caption{Grundy value analysis: (top-left) triangular move Grundy values,
(top-right) perfect square move Grundy values,
(bottom-left) Grundy value distributions,
(bottom-right) $W(N)$ for small $N$.}
\end{figure}

\section{Answer}

\[
\boxed{3996390106631}
\]

\end{document}
